% answers_short.tex —— 1989 年·广东卷·理科 · 简短答案速查
% 由 tools/make_short_answers.py 自 answers.tex 的 \gksolhead/\gkanswer 提取，
% 并与 1989/1989_数学真题_广东理_答案要点.md 逐题比对后生成，勿手改（改源 MD 后重跑）。

\begin{gknotice}
\textbf{使用说明}：本节只给最终答案便于快速核对，完整推导见解析版。答案为编写者独立推导，\textbf{非官方标准答案}。
\end{gknotice}

\gkpart{选择题}
\begin{center}\begin{tabular}{cccccccccccccc}
\toprule
题号 & 1 & 2 & 3 & 4 & 5 & 6 & 7 & 8 & 9 & 10 & 11 & 12 & 13 \\
\midrule
答案 & C & C & A & A & B & D & A & C & C & D & D & B & A \\
\bottomrule
\end{tabular}\end{center}\medskip
\begin{center}\begin{tabular}{ccccccccccccc}
\toprule
题号 & 14 & 15 & 16 & 17 & 18 & 19 & 20 & 21 & 22 & 23 & 24 & 25 \\
\midrule
答案 & D & A & B & B & B & C & A & C & C & D & C & B \\
\bottomrule
\end{tabular}\end{center}

\gkpart{填空题}
\begin{itemize}[leftmargin=2.6em]
  \item[\textbf{一.}] (1) $[0,+\infty)$；(2) $\dfrac{\sqrt2}{2}$；(3) $\dfrac32$；(4) $23$；(5) $6$
\end{itemize}

\gkpart{解答题}
\begin{itemize}[leftmargin=2.6em]
  \item[\textbf{二.}] (1) 见解析版；(2) $5\text{ cm}$
  \item[\textbf{三.}] (1) $(z,\omega)=(-i,-i)$ 或 $(3i,-5i)$；(2) $3\sqrt3$
  \item[\textbf{四.}] (1) $a>2$；(2) 见解析版
  \item[\textbf{五.}] (1) $y^2=x$，其中 $y>0$（不含原点）；(2) 当且仅当 $\theta_2=\dfrac\pi2$ 时，$r_1,r_2,r_3$ 成等比数列
\end{itemize}
